\section{Introduction}
\label{section: introduction}


Quandles, introduced independently by Joyce~\cite{joyce_1982_a_classifying_invariant_of_knots_the_knot_quandle} and Matveev~\cite{matveev_1982_distributive_groupoids_in_knot_theory}, are algebraic structures whose axioms are an abstraction of the Reidemeister moves for knot diagrams.
Besides providing invariants of knots, quandles are closely related to groups.
With a quandle $Q$, one associates several groups, such as the automorphism group $\Aut(Q)$, the inner automorphism group $\Inn(Q)$, the transvection group $\Trans(Q)$, and the adjoint group $\As(Q)$.
Conversely, every group $G$ gives rise to a quandle $\Conj(G)$, called the \emph{conjugacy quandle} of $G$, with underlying set $G$ and operation given by conjugation in $G$.
The functor $\Conj$ is faithful and admits a left adjoint, given by $Q\mapsto \Adj(Q)$.
Thus, quandles provide a broader setting in which groups can be studied through their conjugation structures.

Several of the groups associated with a quandle $Q$ reflect the structure of $Q$ through their natural action on $Q$. Accordingly, one way to study a quandle $Q$ is to study the properties of these groups, or those of their actions on $Q$. For example, a quandle $Q$ is called \emph{homogeneous} if the action $\Aut(Q)\curvearrowright Q$ is transitive, and \emph{connected} if the action $\Inn(Q)\curvearrowright Q$ is transitive. 
Likewise, the \emph{mediality} of a quandle $Q$ is equivalent to the commutativity of $\Trans(Q)$.


In his paper~\cite{darne_2026_nilpotent_quandles}, Darn\'e introduced a notion of nilpotency for quandles, calling a quandle $Q$ \emph{nilpotent} if its inner automorphism group $\Inn(Q)$ is nilpotent.
The resulting class of quandles had appeared earlier in the literature under the terminology of reductive quandles; the use of the word nilpotency is justified by the fact that these quandles behave in close parallel to nilpotent groups.
One manifestation of this parallel is their characterization as the quandles admitting a finite tower of quandle coverings.


Recall that a group $G$ is \emph{nilpotent} if its lower central series becomes trivial after finitely many steps, or equivalently, if $G$ can be obtained from the trivial group by a finite sequence of central group extensions.
The quandle-theoretic counterpart of a central extension is a covering homomorphism, introduced by Eisermann~\cite{eisermann_2014_quandle_coverings_and_their_galois_correspondence} in the context of the covering theory of quandles; coverings of conjugation quandles arising from group extensions correspond to central extensions of groups (see for instance \cite[Proposition 5.2]{preprint_yuki_tomoki_2026_relativization_of_symmetries_on_quandles}).
As an analogue of the above characterization of group nilpotency, Darn\'e showed in \cite[Theorem 2.13]{darne_2026_nilpotent_quandles} that a non-empty quandle $Q$ is nilpotent if and only if the canonical map from $Q$ to the one-element quandle is expressible as a finite composite of covering homomorphisms.

In the present paper, we extend this notion of nilpotency and develop its relative version for surjective homomorphisms of quandles.

In the previous work \cite{preprint_yuki_tomoki_2026_relativization_of_symmetries_on_quandles}, the authors defined the \emph{relative inner automorphism group}
\[ \relInn(f) = \Ker(f_*\colon \Inn(P) \twoheadrightarrow \Inn(Q)) \]
of a surjective quandle homomorphism $f\colon P \twoheadrightarrow Q$, where $f_*$ denotes the induced surjective group homomorphism.
There, a surjective homomorphism $f$ is called \emph{connected} if $\relInn(f)$ acts transitively on every fiber of $f$. 
This is a relativization of the connectedness of quandles, in the sense that a quandle $Q$ is connected if and only if the unique homomorphism $Q\to\{\ast\}$ into the terminal quandle is connected.


% Relativization is not merely a formal generalization.
% なんかいい感じのrelativizeする理由を書く？



Along this line, we define nilpotency for quandle homomorphisms.
It should be noted that the nilpotency of the relative inner automorphism group alone is not sufficient: we must also take into account the conjugation action of the whole inner automorphism group. We therefore make use of the notion of relative nilpotency for a normal subgroup; see \cref{appendix: relative nilpotency for groups} for its general account.


\subsection{Results}
\label{subsection: results in introduction}

We first fix notation. For a normal subgroup $N \trianglelefteq G$, the \emph{$G$-relative lower central series} of $N$ is defined inductively by $\relLCS_0(G,N)=N$ and $\relLCS_{i+1}(G,N)=[G,\relLCS_i(G,N)]$. If $\relLCS_c(G,N)=1$ for some $c\geq 0$, then $N$ is called \emph{$G$-nilpotent}.


Let $f \colon P\twoheadrightarrow Q$ be a surjective quandle homomorphism.
By the surjectivity, there is a unique surjective group homomorphism $f_*\colon \Inn(P)\twoheadrightarrow \Inn(Q)$ satisfying $f_*(s_x)=s_{f(x)}$ for all $x\in P$.
The \emph{relative inner automorphism group} $\relInn(f)$ is defined as $\relInn(f)=\Ker(f_*)$, which is a normal subgroup of $\Inn(P)$.
Then, we call $f$ \emph{nilpotent} if $\relInn(f)$ is nilpotent relative to its ambient group $\Inn(P)$, namely if $\relLCS_{c}(\Inn(P),\relInn(f))=1$ for some $c\geq0$. Thus, relative nilpotency depends not only on the abstract group $\relInn(f)$, but also on the conjugation action of $\Inn(P)$ on it.
For a non-empty quandle $P$ and the unique surjective homomorphism $P\twoheadrightarrow\{\ast\}$ into the terminal quandle, this recovers Darn\'e's nilpotency of $P$.

Our first result generalizes to this relative setting the characterization of nilpotent quandles as those obtained as iterated coverings of the terminal quandle.

\begin{introthm}[{\cref{theorem: nilpotent homomorphism and covering rigid tower}}]
\label[introthm]{introtheorem: nilpotent homomorphism and covering rigid tower}
    Let $f\colon P\twoheadrightarrow Q$ be a surjective quandle homomorphism and let $c\geq0$. Then the following conditions are equivalent.
    \begin{enumerate}
        \item\label{introitem:c-nilpotent} $f$ is $c$-nilpotent.
        \item\label{introitem:factorization_into_c_covering_and_1_rigid} $f$ admits a factorization
        \[
        \begin{tikzcd}
            P=P_{c}
            \arrow[two heads]{r}{p_c}
            &
            P_{c-1}
            \arrow[two heads]{r}{p_{c-1}}
            &
            \cdots
            \arrow[two heads]{r}
            &
            P_{1}
            \arrow[two heads]{r}{p_1}
            &
            P_{0}
            \arrow[two heads]{r}{h}
            &
            Q
        \end{tikzcd}
        \]
        in which the first $c$ homomorphisms $p_c,\dots,p_1$ are coverings and the last homomorphism $h$ is rigid (i.e.~$\relInn(h)=1$).
    \end{enumerate}
\end{introthm}

Since rigid homomorphisms are coverings, \cref{introtheorem: nilpotent homomorphism and covering rigid tower} says that the nilpotent homomorphisms are precisely the finite composites of coverings.

In \cite{preprint_yuki_tomoki_2026_relativization_of_symmetries_on_quandles}, the authors also introduced the \emph{relative transvection group} of $f$,
\[ \relTrans(f) = \left\langle s_xs_y^{-1} \mid x,y \in P \text{ such that }f(x)=f(y) \right\rangle, \]
which is a normal subgroup of $\Inn(P)$ as well. One may use this subgroup to define nilpotency, but the two notions turn out to coincide, because the inclusion $[\Inn(P),\relInn(f)] \subseteq \relTrans(f)$ holds (see \cref{lem:commutator_inn_rel_trans_rel}).
Furthermore, a third group-theoretic description of the nilpotency condition is deduced from the \emph{relative adjoint group} of $f$, which is defined as
\[ \relAs(f)=\Ker\bigl(\As(f)\colon\As(P)\twoheadrightarrow\As(Q)\bigr). \]
Using the fact that the surjective group homomorphism $\rho_P \colon \Adj(P) \twoheadrightarrow\Inn(P)$ is a central extension and sends $\relAs(f)$ onto $\relTrans(f)$, we prove that the relative nilpotency of $\relAs(f)$ is equivalent to that of $\relTrans(f)$.
 
Summarizing, we establish the following theorem.

\begin{introthm}[{\cref{theorem: main structure theorem for nilpotent homomorphisms}}]
\label[introthm]{introtheorem: main structure theorem for nilpotent homomorphisms}
    Let $f\colon P\twoheadrightarrow Q$ be a surjective quandle
    homomorphism. Then the following conditions are equivalent.
    \begin{enumerate}[label={(\roman*)}]
        \item\label{introitem: f is nilp} $f$ is nilpotent, that is, $\relInn(f)$ is $\Inn(P)$-nilpotent.
        \item\label{introitem: reltrans is nilp} $\relTrans(f)$ is $\Inn(P)$-nilpotent.
        \item\label{introitem: relas is nilp} $\relAs(f)$ is $\As(P)$-nilpotent.
        \item\label{introitem: composition of coverings} $f$ is a finite composition of covering homomorphisms.
    \end{enumerate}
\end{introthm}

Considering the case with $Q=\{\ast\}$, we also obtain equivalent descriptions of nilpotency for a non-empty quandle (\cref{corollary: characterization_of_nilpitent_quandles}). In particular, a non-empty quandle $P$ is nilpotent (i.e., $\Inn(P)$ is nilpotent) if and only if $\Trans(P)$ is $\Inn(P)$-nilpotent.
This is a novel equivalent condition for nilpotency, which is not included in \cite[Corollary~2.5 and Theorem~2.13]{darne_2026_nilpotent_quandles}.

It is worth mentioning that Bonatto and Stanovsk\'{y}~\cite{bonatto_stanovsky_2021_commutator_theory_for_racks_and_quandles} also study nilpotency for quandles from the universal-algebraic viewpoint. According to \cite[Theorem 1.2]{bonatto_stanovsky_2021_commutator_theory_for_racks_and_quandles}, the nilpotency of a quandle $P$ in their sense is equivalent to that of $\Trans(P)$. Hence, our notion of nilpotency is stronger than theirs, as noted by Darn\'{e} in the introduction of \cite{darne_2026_nilpotent_quandles}.




%% section 3

Next, we extend the relative lower central series to transfinite ordinals.
Even when it does not reach $1$ after finitely many steps, the series may still do so at an infinite stage; for groups, this weaker condition is classical and known as hypocentrality.
So, we call $f$ a \emph{hypocentral} homomorphism if the transfinite $\Inn(P)$-relative lower central series of $\relInn(f)$ reaches $1$ at some, possibly infinite, ordinal stage. Every nilpotent homomorphism is hypocentral, but the converse fails in general.
We obtain the following group-theoretic characterization of hypocentrality.


\begin{introthm}[\cref{theorem: characterizations of hypocentral homomorphisms}]
\label[introthm]{introtheorem: characterizations of hypocentral homomorphisms}
Let $f\colon P\twoheadrightarrow Q$ be a surjective quandle homomorphism. Then, the following conditions are equivalent.
\begin{enumerate}[label={(\roman*)}]
    \item\label{introitem: f is hypocentral}
    The homomorphism $f$ is hypocentral, that is, $\relInn(f)$ is $\Inn(P)$-hypocentral.
    \item\label{introitem: relative transvection is hypocentral}
    The group $\relTrans(f)$ is $\Inn(P)$-hypocentral.
    \item\label{introitem: relative adjoint group is hypocentral}
    The group $\relAs(f)$ is $\As(P)$-hypocentral.
    \item\label{introitem: no perfect subgroup in relative inner group}
    The group $\relInn(f)$ contains no non-trivial $\Inn(P)$-perfect
    normal subgroup of $\Inn(P)$.
\end{enumerate}
\end{introthm}


Passing from nilpotency to hypocentrality clarifies the relationship with the notion of strongly connected homomorphisms.
In \cite[Definition 3.6]{preprint_yuki_tomoki_2026_relativization_of_symmetries_on_quandles}, a surjective homomorphism $f\colon P\twoheadrightarrow Q$ is called \emph{strongly connected} if the relative transvection group $\relTrans(f)$ acts transitively on every fiber of $f$.
The authors observed in \cite[Proposition 3.21]{preprint_yuki_tomoki_2026_relativization_of_symmetries_on_quandles} that the class of strongly connected homomorphisms is left orthogonal to the class of pre-covering homomorphisms. Here, a \emph{pre-covering} homomorphism is a not necessarily surjective homomorphism $f$ such that $f(x)=f(y)$ implies $s_x=s_y$.
We extend this observation by showing that hypocentral homomorphisms admit a category-theoretic characterization as the class of morphisms right orthogonal to strongly connected homomorphisms.


\begin{introthm}[\cref{theorem: transfinite covering characterization of hypocentral homomorphisms}]
\label[introthm]{introtheorem: transfinite covering characterization of hypocentral homomorphisms}
    Let $f\colon P\twoheadrightarrow Q$ be a surjective quandle homomorphism. Then, the following conditions are equivalent.
\begin{enumerate}[label={(\roman*)}]
    \item\label{introitem: homomorphism is hypocentral}
    $f$ is hypocentral.
    \item\label{introitem: transfinite cocomposite of covering condition}
    $f$ is a transfinite cocomposite of pre-covering homomorphisms.
    \item\label{introitem: hypocentral right orthogonal}
    $f$ is right orthogonal to every strongly connected homomorphism.
\end{enumerate}
\end{introthm}
The use of pre-covering homomorphisms is essential: a hypocentral homomorphism need not be a transfinite cocomposite of coverings (see \cref{example: hypocentral but not nilpotent Alexander quandle}).

Even and Gran~\cite[Proposition 3.2]{even_gran_2014_on_factorization_systems_for_surjective_quandle_homomorphisms} previously showed that the classes of connected and rigid homomorphisms form an orthogonal factorization system relative to surjective homomorphisms (see also \cite[Proposition 2.24]{preprint_yuki_tomoki_2026_relativization_of_symmetries_on_quandles}). 
We establish the analogous statement for strongly connected homomorphisms.

\begin{introthm}[Strongly connected--hypocentral factorization system; \cref{theorem: universal strongly connected-hypocentral factorization,theorem: strongly connected-hypocentral factorization system}]
\label[introthm]{introtheorem: strongly connected-hypocentral factorization system}
    Every surjective quandle homomorphism $f$ admits a factorization $f=\hypFac{f}\circ \hypMap{f}$ into a strongly connected homomorphism $\hypMap{f}$ followed by a hypocentral homomorphism $\hypFac{f}$.
    Thus, the pair $(\{\mathrm{strongly\>connected}\}, \{\mathrm{hypocentral}\})$ of classes of surjections is an orthogonal factorization system for the surjections in $\Quandle$.
\end{introthm}


From \cite[Proposition 3.32]{preprint_yuki_tomoki_2026_relativization_of_symmetries_on_quandles}, strongly connected homomorphisms are closed under pullbacks along surjective homomorphisms. Therefore, this factorization system for the surjections is in fact \textit{stable} in the sense of \cite[\href{https://ncatlab.org/nlab/show/stable+factorization+system}{stable factorization system}]{preprint_nlab_2008_nlab}.

%We further show that the hypocentral factor $\hypFac{f}$ is nilpotent precisely when the relative transvection series of $f$ stabilizes at a finite stage. In particular, hypocentrality and nilpotency agree for every surjection between quandles with finite inner automorphism groups, so that the classes of strongly connected and nilpotent homomorphisms already form a stable orthogonal factorization system on this subcategory, for instance on that of finite quandles. This is no longer true in general: for the Alexander quandle associated with $-\id$ on $\mathbb{Z}$, the terminal homomorphism is hypocentral but not nilpotent.


%%% section 4

In the third part of the paper, we introduce relative analogues of reduced quandles. Reduced quandles, also called quasi-trivial quandles in part of the literature, arise in the study of link-homotopy \cite{hughes_2011_link_homotopy_invariant_quandles,inoue_2013_quasitriviality_of_quandles_for_linkhomotopy,elhamdadi_liu_nelson_2018_quasitrivial_quandles_and_biquandles_cocycle_enhancements_and_linkhomotopy_of_pretzel_links} and are studied systematically in \cite[Section~7]{darne_2026_nilpotent_quandles}.
We study their relativization, called reduced homomorphisms, together with the slightly more general operator-reduced homomorphisms; these two notions lie between the first and the second nilpotency.

A surjective homomorphism $f$ is called \emph{reduced} if $s_{\varphi(x)}(x)=x$, and \emph{operator-reduced} if $s_xs_{\varphi(x)}=s_{\varphi(x)}s_x$, for every $x\in P$ and every $\varphi\in\relInn(f)$.
As we will see in \cref{proposition: commutator interpretation of reducedness}, $1$-nilpotent implies reduced, reduced implies operator-reduced, and $2$-nilpotent implies operator-reduced.
We then investigate when an operator-reduced homomorphism is nilpotent, and establish the following ascent theorem.

\begin{introthm}[Nilpotency ascent theorem;~\cref{theorem: finite operator reduced ascent}]
\label[introthm]{introtheorem: finite operator reduced ascent}
    Let $f\colon P\twoheadrightarrow Q$ be an operator-reduced quandle homomorphism such that $\relInn(f)$ is a finite group. If $Q$ is nilpotent, then both $f$ and $P$ are nilpotent.
\end{introthm}
In Darn\'e's absolute setting \cite[Proposition 7.13]{darne_2026_nilpotent_quandles}, the target is the one-point quandle and hence automatically nilpotent. 
The relative viewpoint reveals the importance of this condition: even a connected reduced homomorphism with trivial fibers and a finite relative inner automorphism group need not be nilpotent (see \cref{example: finite connected reduced homomorphism not nilpotent}).

Under some technical assumptions on the commutativity of the conjugation action on abelian sections of $\relInn(f)$, we also prove another criterion for the nilpotency of an operator-reduced homomorphism (\cref{theorem: commutative action criterion}).
In contrast with \cref{introtheorem: finite operator reduced ascent}, it requires no finiteness assumption on $\relInn(f)$.
Applying this criterion, we obtain the following corollary, which covers a number of natural situations.

\begin{introcor}[\cref{lemma: abelian images of inner automorphism groups,corollary: easy consequences of comm action criterion}]
\label[introcor]{introcorollary: easy consequences of comm action criterion}
    Let $f\colon P\twoheadrightarrow Q$ be an operator-reduced quandle homomorphism. 
    Suppose that $P$ satisfies either the condition that $\Inn(P)$ is generated by finitely many symmetry maps or the condition that $P$ has finitely many connected components.
    If $\relInn(f)$ is nilpotent and $\Inn(Q)$ is abelian, then $f$ is nilpotent.
\end{introcor}


%The nilpotency of the target is essential in \cref{introtheorem: finite operator reduced ascent}. We exhibit a finite connected reduced homomorphism with trivial fibers and abelian relative inner automorphism group which is nevertheless not nilpotent. Hence, neither the triviality of the fibers nor the abstract nilpotency of $\relInn(f)$ controls relative nilpotency, as long as no condition is imposed on the action coming from the base.




\begin{Out}
\cref{section: preliminaries on quandles} collects the preliminaries on quandles used throughout the paper. We also recall the relative viewpoint on surjective quandle homomorphisms introduced in \cite{preprint_yuki_tomoki_2026_relativization_of_symmetries_on_quandles}.
 
In \cref{section: nilpotent homomorphisms}, we introduce the notion of nilpotency for surjective quandle homomorphisms and develop its basic theory. After recalling Darn\'e's nilpotent quandles, we define nilpotent homomorphisms through the relative lower central series of $\relInn(f)$ and prove \cref{introtheorem: nilpotent homomorphism and covering rigid tower}.
%, together with the behaviour of nilpotency under composition, products, pullbacks and passage to fibers, as well as the construction of the universal $c$-nilpotent truncation.
We then give the descriptions of nilpotency in terms of the relative transvection group and of the relative adjoint group, which yield \cref{introtheorem: main structure theorem for nilpotent homomorphisms}.
%; along the way, the relative transvection group is shown to compute the covering length.
%The section closes with a comparison of the three resulting numerical invariants and with explicit computations for homomorphisms of Alexander quandles, which show that the bounds obtained are sharp.
 
\cref{section: hypocentral homomorphisms} extends the relative lower central series to transfinite ordinals and introduces hypocentral homomorphisms. We first establish their basic properties and prove \cref{introtheorem: characterizations of hypocentral homomorphisms}. We then construct the strongly connected--hypocentral factorization and prove \cref{introtheorem: transfinite covering characterization of hypocentral homomorphisms,introtheorem: strongly connected-hypocentral factorization system}.
%In the last part, we determine when the hypocentral factor is nilpotent; in particular, the two classes coincide under a descending chain condition, which applies to all finite quandles, whereas an Alexander quandle over $\mathbb{Z}$ provides a hypocentral homomorphism that is not nilpotent.
 
\cref{section: relative reduced homomorphisms} is devoted to reduced and operator-reduced homomorphisms. We first define these two classes and locate them between $1$-nilpotency and $2$-nilpotency.
%, and establish their closure properties under composition, products and pullbacks.
We next construct the universal reduced and operator-reduced factors of a surjective homomorphism. 
Finally, we prove \cref{introtheorem: finite operator reduced ascent} and the commutative action criterion (\cref{theorem: commutative action criterion}), from which \cref{introcorollary: easy consequences of comm action criterion} is deduced, and we conclude with the case of a cyclic relative inner automorphism group.
 
\cref{appendix: relative nilpotency for groups} collects the purely group-theoretic material on the relative nilpotency of a normal subgroup, or equivalently on nilpotent surjections of groups.
\end{Out}





\begin{NoCon}
Groups act on the left. For elements $x,y$ of a group, we use the commutator convention $[x,y]=xyx^{-1}y^{-1}$. For subgroups $A$ and $B$, the notation $[A,B]$ denotes the subgroup generated by the commutators $[a,b]$ with $a\in A$ and $b\in B$.
\end{NoCon}

\begin{AI}
The authors used GPT-5.6 Sol and GPT-6 Astra (OpenAI), Gemini 3.1 Pro (Google), and Claude Opus 5 (Anthropic) throughout the preparation of this work, particularly for the proof of \cref{theorem: finite operator reduced ascent}, the construction of \cref{example: hypocentral but not nilpotent Alexander quandle}, and for English proofreading. 
All arguments and proofs were independently verified by the authors, who take full responsibility for the contents of this paper.
\end{AI}

\begin{Ack}
This work was supported by JSPS KAKENHI Grant Number JP25K23333 and JP26KJ0276.
The first author acknowledges access to Claude through the Claude Team plan for scientists and the second author acknowledges support from OpenAI through the ChatGPT for Academic Researchers program.
\end{Ack}
